%lesson 2, 05/03/2026

\subsection{Autovettori e autovalori}

Dato un operatore autoaggiunto $ A $, si consideri ora l'equazione agli autovalori
$$ A\ket{a} = a\ket{a} $$ 

Prima di dare la definizione di valori regolari e autovalori di $ A $, si consideri la seguente
quantità, sempre definita:
$$ C := \inf\limits_{\ket{\psi}\in\mathcal{D}_A}\frac{||(A-a1)\ket{\psi}||}{||\ket{\psi}||}\geq 0$$ 
\textbf{Def.} $a\in\mathbb{R}$ si dice punto regolare di $ A $ se $ C > 0 $.

\textbf{Def.} $a\in\mathbb{R}$ si dice autovalore di $ A $ se $C=0$

Poiché $ C $ è un estremo inferiore si distinguono due casi:
\begin{itemize}
    \item $ C $ è anche il minimo: in questo caso $ a $ è un \textbf{autovalore proprio} di $ A $
	ed $ \exists \ket{a}\in\mathcal{D}_A $ tale che $ A\ket{a} = a\ket{a}$, detto
	\textbf{autovettore proprio}.
    \item $ C $ non è il minimo: in questo caso $ a $ è un \textbf{autovalore generalizzato} di 
	$ A $ e l'autovettore corrispondente $ \ket{a}\notin\mathcal{D}_A $ è detto
	\textbf{autovettore generalizzato} . 
\end{itemize}

L'insieme degli	autovalori propri è detto \textbf{spettro discreto di A} \textit{Spd}(A).

L'insieme degli autovalori generalizzati è detto \textbf{spettro continuo di A} \textit{Spc}(A).

L'importanza degli spettri di un operatore è evidenziata dal seguente risultato, detto appunto
\textit{teorema spettrale}.

\noindent\fbox{%
    \parbox{\textwidth}{%
    \textbf{Teorema spettrale}

    Dato un operatore $ A $ autoaggiunto, le soluzioni dell'equazione $ A\ket{a} = a\ket{a} $ sono
    un insieme \textbf{completo}.

    $ \forall f \in S $ si definiscono 
\begin{itemize}
    \item I coefficienti $ f_i := \bra{a_i}\ket{f} $, con $ \ket{a_i} $ autovettori propri
	corrispondenti agli autovalori

	$ a_i \in$ \textit{Spd}(A);
    \item La funzione $ f(a):=\bra{a}\ket{f} $, con $ \ket{a} $ autovettori generalizzati 
	corrispondenti agli autovalori $ a\in $ \textit{Spc}(A).
\end{itemize}

Si può scrivere $ \forall f,g\in S $
\begin{flalign*}
 & \bra{f}\ket{g} = \sum_{i}^{}f^*_ig_i + \int_{Spd(A)}^{}daf^*(a)g(a)
    := \sumint_{Sp(A)}daf^*(a)g(a)\text{ , prodotto scalare}&&\\
& \ket{f} = \sumint_{Sp(A)}daf(a)\ket{a}\text{ , rappresentazione di $\ket{f}$}&&\\
 & \sumint_{Sp(A)}\ket{a}\bra{a} = 1 \text{ , relazione di completezza dello spettro} &&
\end{flalign*}


    }%
}

